\chapter{Lesson 3: Weather and Seasons}

Talking about the weather is the universal small-talk topic, and it was a natural one in a
Waterloo winter term. This lesson adds descriptive adjectives, the two clause connectors
\ko{\til 고} (and) and \ko{\til 지만} (but), and a first look at the \ko{ㅂ}-irregular
adjectives, which include the two most important weather words, \ko{덥다} (hot) and
\ko{춥다} (cold).

\section{Vocabulary}

\begin{table}[H]
\centering
\begin{tabular}{lll}
\toprule
Korean & RR & English \\
\midrule
\ko{날씨} & nalssi & weather \\
\ko{봄 / 여름 / 가을 / 겨울} & bom / yeoreum / gaeul / gyeoul & spring / summer / autumn / winter \\
\ko{계절} & gyejeol & season \\
\ko{덥다} & deopda & to be hot (weather) \\
\ko{춥다} & chupda & to be cold (weather) \\
\ko{따뜻하다} & ttatteutada & to be warm \\
\ko{시원하다} & siwonhada & to be cool, refreshing \\
\ko{맑다} & makda & to be clear, sunny \\
\ko{흐리다} & heurida & to be cloudy \\
\ko{비가 오다} & biga oda & to rain (lit.\ rain comes) \\
\ko{눈이 오다} & nuni oda & to snow \\
\ko{바람이 불다} & barami bulda & to be windy (lit.\ wind blows) \\
\ko{습하다} & seupada & to be humid \\
\ko{기온} & gion & temperature \\
\ko{도} & do & degree(s) \\
\ko{영하} & yeongha & below zero \\
\ko{일기예보} & ilgiyebo & weather forecast \\
\ko{벚꽃} & beotkkot & cherry blossoms \\
\ko{단풍} & danpung & autumn leaves \\
\ko{장마} & jangma & monsoon (rainy) season \\
\bottomrule
\end{tabular}
\end{table}

\begin{remark}
\ko{덥다/춥다} describe the air temperature or how your body feels. For objects use
\ko{뜨겁다} tteugeopda (hot to the touch) and \ko{차갑다} chagapda (cold to the touch):
\ko{커피가 뜨거워요} (the coffee is hot). For water, food and drinks that are pleasantly
cool, use \ko{시원하다}.
\end{remark}

\subsection{The seasons in Korea}

\begin{center}
\begin{tikzpicture}[every node/.style={font=\small}]
  \fill[pink!40] (0,0) -- (90:2.4) arc (90:0:2.4) -- cycle;
  \fill[orange!35] (0,0) -- (0:2.4) arc (0:-90:2.4) -- cycle;
  \fill[brown!35] (0,0) -- (-90:2.4) arc (-90:-180:2.4) -- cycle;
  \fill[blue!15] (0,0) -- (180:2.4) arc (180:90:2.4) -- cycle;
  \draw[thick] (0,0) circle (2.4);
  \draw (0,-2.4) -- (0,2.4); \draw (-2.4,0) -- (2.4,0);
  \node[align=center] at (45:1.4) {\ko{봄}\\ \scriptsize 3--5\ko{월}};
  \node[align=center] at (-45:1.4) {\ko{여름}\\ \scriptsize 6--8\ko{월}};
  \node[align=center] at (-135:1.4) {\ko{가을}\\ \scriptsize 9--11\ko{월}};
  \node[align=center] at (135:1.4) {\ko{겨울}\\ \scriptsize 12--2\ko{월}};
  \node[anchor=west, font=\scriptsize] at (2.6,1.5) {\ko{따뜻해요}, \ko{벚꽃}, \ko{황사} (yellow dust)};
  \node[anchor=west, font=\scriptsize] at (2.6,-1.5) {\ko{더워요}, \ko{습해요}, \ko{장마} (late June--July)};
  \node[anchor=east, font=\scriptsize] at (-2.6,-1.5) {\ko{시원해요}, \ko{맑아요}, \ko{단풍}};
  \node[anchor=east, font=\scriptsize] at (-2.6,1.5) {\ko{추워요}, \ko{눈이 와요}, \ko{영하}};
\end{tikzpicture}
\end{center}

\section{Asking about and describing the weather}

\begin{example}[Weather questions and answers]
The two questions to know are \ko{날씨가 어때요?} Nalssiga eottaeyo? (How is the
weather?) and \ko{몇 도예요?} Myeot doyeyo? (What's the temperature?).
\begin{itemize}
  \item \ko{오늘 날씨가 어때요? --- 맑아요. 그런데 좀 추워요.} Oneul nalssiga eottaeyo?
        --- Malgayo. Geureonde jom chuwoyo. (How's the weather today? --- It's clear, but
        a bit cold.)
  \item \ko{지금 영하 십오 도예요.} Jigeum yeongha sibo doyeyo. (It's minus fifteen
        degrees right now.) Temperatures use Sino-Korean numbers.
  \item \ko{어제 눈이 많이 왔어요.} Eoje nuni mani wasseoyo. (It snowed a lot
        yesterday.)
  \item \ko{일기예보에서 내일 비가 온다고 했어요.} Ilgiyeboeseo naeil biga ondago
        haesseoyo. (The forecast said it'll rain tomorrow.) --- we were only taught this
        as a fixed phrase.
\end{itemize}
\end{example}

\section{Grammar 1: \ko{\til 고} (and)}

\ko{\til 고} attaches directly to a verb or adjective stem and joins two clauses. It has
two uses:
\begin{enumerate}
  \item \textbf{Listing} facts (``A and B''), usually with adjectives:
        \ko{날씨가 맑고 따뜻해요.} Nalssiga malgo ttatteutaeyo. (The weather is clear and
        warm.)
  \item \textbf{Sequence} of actions (``A and then B''):
        \ko{밥을 먹고 학교에 가요.} Babeul meokgo hakgyoe gayo. (I eat and then go to
        school.)
\end{enumerate}

\begin{center}
\begin{tikzpicture}[every node/.style={font=\small},
    c/.style={draw, rounded corners, minimum height=0.8cm, fill=gray!8}]
  \node[c] (a) at (0,0) {\ko{숙제를 하}};
  \node[c, fill=orange!15, right=0pt of a] (g) {\ko{고}};
  \node[c] (b) at (4.2,0) {\ko{친구를 만났어요}};
  \node[font=\scriptsize, below=0.25cm of a] {no tense here};
  \node[font=\scriptsize, below=0.25cm of b] {tense marked once, at the end};
  \draw[-{Stealth}, thick] (g.east) -- (b.west);
  \node[font=\scriptsize] at (2.3,0.85) {sukjereul hago chingureul mannasseoyo. ``I did homework and (then) met a friend.''};
\end{tikzpicture}
\end{center}

\begin{important}[Tense with \ko{\til 고}]
In a \emph{sequence}, the past tense appears only on the final verb: \ko{어제 숙제를 하고
잤어요} (yesterday I did my homework and slept), not \ko{*했고 잤어요}. For a plain
\emph{list} of two separate past facts, the first verb may also take the past:
\ko{어제는 바람이 불었고 눈도 왔어요} (yesterday it was windy, and it snowed too).
\end{important}

\section{Grammar 2: \ko{\til 지만} (but)}

\ko{\til 지만} attaches to the stem and contrasts two clauses, like \ko{그렇지만} at the
start of a sentence. Unlike sequential \ko{\til 고}, it \emph{does} take the past tense
on the first clause when the first event is past: \ko{\til 았/었지만}.

\begin{example}[Contrast]
Stem + \ko{지만}:
\begin{itemize}
  \item \ko{여름은 덥지만 저는 여름이 좋아요.} Yeoreumeun deopjiman jeoneun yeoreumi
        joayo. (Summer is hot, but I like summer.)
  \item \ko{한국어는 어렵지만 재미있어요.} Hangugeoneun eoryeopjiman jaemiisseoyo. (Korean
        is difficult but fun.)
  \item \ko{어제는 비가 왔지만 오늘은 맑아요.} Eojeneun biga watjiman oneureun malgayo.
        (It rained yesterday, but today is sunny.)
  \item \ko{학생이지만 일도 해요.} Haksaengijiman ildo haeyo. (I'm a student, but I
        work too.) --- \ko{이다} + \ko{지만} $\to$ \ko{이지만}.
\end{itemize}
\end{example}

\section{Grammar 3: \ko{ㅂ}-irregular adjectives (first look)}

Many common adjectives whose stems end in \ko{ㅂ} are irregular: before a vowel-initial
ending (\ko{\til 아요/어요}, \ko{\til 았/었어요}) the \ko{ㅂ} turns into \ko{우}, which
then combines with \ko{어} to give \ko{워}. Before a consonant-initial ending
(\ko{\til 고}, \ko{\til 지만}) nothing happens.

\begin{table}[H]
\centering
\begin{tabular}{llllll}
\toprule
Dictionary & \ko{\til 어요} & \ko{\til 었어요} & \ko{\til 고} & \ko{\til 지만} & Meaning \\
\midrule
\ko{덥다} & \ko{더워요} & \ko{더웠어요} & \ko{덥고} & \ko{덥지만} & hot \\
\ko{춥다} & \ko{추워요} & \ko{추웠어요} & \ko{춥고} & \ko{춥지만} & cold \\
\ko{쉽다} & \ko{쉬워요} & \ko{쉬웠어요} & \ko{쉽고} & \ko{쉽지만} & easy \\
\ko{어렵다} & \ko{어려워요} & \ko{어려웠어요} & \ko{어렵고} & \ko{어렵지만} & difficult \\
\ko{맵다} & \ko{매워요} & \ko{매웠어요} & \ko{맵고} & \ko{맵지만} & spicy \\
\ko{귀엽다} & \ko{귀여워요} & \ko{귀여웠어요} & \ko{귀엽고} & \ko{귀엽지만} & cute \\
\bottomrule
\end{tabular}
\end{table}

\begin{remark}
\ko{입다} (to wear), \ko{잡다} (to catch) and \ko{좁다} (to be narrow) are
\emph{regular}: \ko{입어요}, \ko{잡아요}, \ko{좁아요}. There is no way to tell from the
spelling; you learn which is which. The complete set of irregular verb classes is in the
last chapter.
\end{remark}

\section{Grammar 4: \ko{\til 네요} (noticing something)}

The ending \ko{\til 네요} attaches to the stem and expresses mild surprise or a reaction
to something you have just noticed, a bit like ``Oh, it's ...!''. It is very common in
weather talk.

\begin{example}[Reacting]
On looking out the window:
\begin{itemize}
  \item \ko{와, 눈이 오네요!} Wa, nuni oneyo! (Wow, it's snowing!)
  \item \ko{오늘 정말 춥네요.} Oneul jeongmal chumneyo. (It's really cold today.) Note
        nasalization: \ko{춥네요} [\ko{춤네요}].
  \item \ko{한국어를 잘하시네요!} Hangugeoreul jalhasineyo! (You speak Korean well!)
\end{itemize}
\end{example}

\section{Reading: a postcard}

\begin{quote}
\ko{토론토는 지금 겨울이에요. 날씨가 아주 춥고 눈이 많이 와요. 어제는 영하 이십 도였어요!
춥지만 눈이 예뻐요. 서울 날씨는 어때요? 서울도 추워요?}
\end{quote}
Toronto-neun jigeum gyeourieyo. Nalssiga aju chupgo nuni mani wayo. Eojeneun yeongha isip
doyeosseoyo! Chupjiman nuni yeppeoyo. Seoul nalssineun eottaeyo? Seouldo chuwoyo?
(It's winter in Toronto now. The weather is very cold and it snows a lot. Yesterday it
was minus twenty! It's cold, but the snow is pretty. How's the weather in Seoul? Is it
cold in Seoul too?)

\begin{aside}
\textbf{Weather and food.} Koreans eat hot \ko{삼계탕} samgyetang (ginseng chicken soup) on
the hottest days of summer, following the saying \ko{이열치열} iyeolchiyeol, ``fight heat
with heat''. Rainy days call for \ko{파전} pajeon (green-onion pancake) with
\ko{막걸리} makgeolli; the sizzling of the pancake is said to sound like rain. Winter
used to follow a rhythm called \ko{삼한사온} samhansaon, ``three cold days, four warm
days''. In spring, people check the fine-dust (\ko{미세먼지} misemeonji) forecast as
carefully as the temperature.
\end{aside}

\section*{Review summary}
\begin{itemize}
  \item \ko{날씨가 어때요?} --- \ko{맑아요, 흐려요, 비가 와요, 눈이 와요, 바람이 불어요}.
  \item \ko{\til 고}: list (``and'') or sequence (``and then''); tense goes on the last
        verb of a sequence.
  \item \ko{\til 지만}: ``but''; can take past \ko{\til 았/었지만}.
  \item \ko{ㅂ}-irregular: \ko{ㅂ} + vowel ending $\to$ \ko{워}
        (\ko{덥다} $\to$ \ko{더워요}); regular before consonant endings.
  \item \ko{\til 네요}: reacting to something just noticed.
\end{itemize}
